\section{The exact gamma reference and the arithmetic Gram determinants}
\label{sec:theta-gamma-reference}

The norm sequence in (KL.24) is attached to the actual measure
\[
 d\mu(t)=|g(\tfrac12+it)|^2\frac{dt}{2\pi},\qquad g=2\xi.
                                                               \tag{TG.1}
\]
We construct an explicit gamma reference for this measure, prove its
polynomial norms, and retain the full arithmetic multiplier in finite
Gram determinants.  This provides another exact formula for the
original boundary allowance; no asymptotic estimate for an arithmetic
determinant is assumed.

\subsection{The original measure, phase, and packet cancellation}

The defining identity for the completed function gives, on the original
line \(s=1/2+it\),
\[
 \begin{aligned}
 g(\tfrac12+it)&=\Gamma(\tfrac14+it/2)\,a(t),\\
 a(t)&=-\pi^{-1/4}(t^2+\tfrac14)
           e^{-it\log\pi/2}\zeta(\tfrac12+it).
 \end{aligned}                                                   \tag{TG.2}
\]
Indeed \(s(s-1)=-(t^2+1/4)\) and
\(\pi^{-s/2}=\pi^{-1/4}e^{-it\log\pi/2}\).  Both the minus
sign and this phase remain in the multiplier.

For \(\lambda>0\), define
\[
 d\sigma_\lambda(t)=|\Gamma(\lambda+it/2)|^2\frac{dt}{2\pi}.
                                                               \tag{TG.3}
\]
In particular
\[
 d\mu=|a|^2d\sigma_{1/4},\qquad
 d\nu_h=|a_h|^2d\sigma_{1/4},\qquad
 a_h(t)=\frac{g(1/2+it)/h(1/2+it)}{\Gamma(1/4+it/2)}.
                                                               \tag{TG.4}
\]
The last expression defines the values at every zero of the original
full packet \(h\); away from those points it equals
\(a(t)/h(1/2+it)\).  The function \(g/h\) is entire because
the complete order of each selected zero is removed.  Gamma has no
zero and has no pole on the real \(t\) axis.  Thus (TG.4) retains
the actual cancellation at each critical packet centre, without
introducing a pole into the measure.

Multiplication gives surjective Hilbert isometries
\[
 \begin{aligned}
 U_a:L^2(\mu)&\longrightarrow L^2(\sigma_{1/4}),
                  &F&\longmapsto aF,\\
 U_{a_h}:L^2(\nu_h)&\longrightarrow L^2(\sigma_{1/4}),
                  &F&\longmapsto a_hF.
 \end{aligned}                                                    \tag{TG.5}
\]
The integral identity proves the isometry.  The zero set of each
multiplier on the real axis is discrete, since it comes from a
nonzero analytic function.  Given \(H\in L^2(\sigma_{1/4})\),
define \(F=H/a\), respectively \(H/a_h\), off this null set
and give \(F\) arbitrary finite values on it.  Then
\(\int |F|^2d\mu=\int |H|^2d\sigma_{1/4}\), respectively
the packet identity, proving surjectivity and the stated inverse.
These are maps between the indicated Hilbert spaces.  Their images
of polynomial subspaces are exactly \(a\mathbb C[t]\) and
\(a_h\mathbb C[t]\), with the displayed multipliers retained.
Moreover the packet isometry from (KL.23) satisfies the exact diagram
\[
 U_{a_h}\bigl(h(1/2+it)P(t)\bigr)=U_aP(t).
                                                               \tag{TG.6}
\]

\subsection{Gamma norms from a generating integral}

We give the integral derivation to fix all constants independently of
different conventions for Meixner--Pollaczek polynomials.  The beta
integral, with \(r=e^x\), gives
\[
 \int_{\mathbb R}\frac{e^{\lambda x}}{(1+e^x)^{2\lambda}}
                     e^{iux}\,dx
       =\frac{\Gamma(\lambda+iu)\Gamma(\lambda-iu)}
                    {\Gamma(2\lambda)}.
                                                               \tag{TG.7}
\]
For completeness the beta integral follows by applying the change of
variables \(v=r_1+r_2\), \(y=r_1/(r_1+r_2)\) to the product of
the two absolutely convergent gamma integrals; the Jacobian is \(v\).
The conditions \(\operatorname{Re}(\lambda\pm iu)>0\) hold for
real \(u\).  The function on the left before Fourier transformation
and its second derivative are integrable: at each end they are
bounded by a constant times \(e^{-\lambda|x|}\).
Two integrations by parts therefore make its Fourier transform
integrable.  Fourier inversion applies.  One direct justification is
to multiply the transform by \(e^{-\varepsilon u^2}\), apply
Fubini, and use the Gaussian approximate identity; dominated
convergence removes that factor because the transform is integrable.
Using its evenness and then \(t=2u\), we obtain
\[
 \int_{\mathbb R}e^{iyt}\,d\sigma_\lambda(t)
   =\frac{2\Gamma(2\lambda)}{(2\cosh y)^{2\lambda}}
   =2^{1-2\lambda}\Gamma(2\lambda)(\cosh y)^{-2\lambda}.
                                                               \tag{TG.8}
\]
The factor \(2\) in the numerator is the Jacobian of the original
change \(t=2u\).

The following integral proves the required exponential integrability.

The function \((2\cosh(x/2))^{-2\lambda}\) extends holomorphically
to \(|\operatorname{Im}x|<\pi\), using the branch positive on
the real axis.  On each closed narrower strip it is bounded by
\(C e^{-\lambda|\operatorname{Re}x|}\).  In (TG.7), for
\(u>0\), shift the contour to \(\mathbb R+ic\), where
\(0<c<\pi\).  The vertical sides of the truncated rectangle
tend to zero by that bound.  For \(u<0\), shift down instead.
It follows that
\[
 |\Gamma(\lambda+iu)|^2\leq C_{\lambda,c}e^{-c|u|}.
\]
Consequently \(\int e^{v|t|}d\sigma_\lambda(t)<\infty\) for
every \(0<v<\pi/2\), by choosing \(2v<c<\pi\).
This proves holomorphic continuation of both sides of (TG.8) to
\(|\operatorname{Im}y|<\pi/2\), where they agree by the identity
theorem and use the branch inherited from the real axis.

For \(z\) near zero use the logarithms that vanish at \(z=0\) and set
\[
 F_\lambda(z,t)=(1-iz)^{-\lambda+it/2}
                        (1+iz)^{-\lambda-it/2}
      =\sum_{n\geq0}P_n^{(\lambda)}(t/2;\pi/2)z^n.
                                                               \tag{TG.9}
\]
For real \(z\) near zero this is
\((1+z^2)^{-\lambda}\exp(t\arctan z)\).  Apply (TG.8) at
\(y=-i(\arctan z+\arctan w)\), where \(z,w\) are real and
sufficiently small.  The identity
\[
 \cos(\arctan z+\arctan w)
        =\frac{1-zw}{\sqrt{1+z^2}\sqrt{1+w^2}}
\]
then proves
\[
 \int_{\mathbb R}F_\lambda(z,t)F_\lambda(w,t)
                            d\sigma_\lambda(t)
   =2^{1-2\lambda}\Gamma(2\lambda)(1-zw)^{-2\lambda}.
                                                               \tag{TG.10}
\]
All derivatives in \(z,w\) at zero can be passed through the
integral: on a smaller closed bidisc their integrands are bounded
by a polynomial in \(|t|\) times \(e^{c|t|}\), with
\(c<\pi/2\).  Thus comparison of the Taylor coefficients in
(TG.10) is justified, not formal.

Define the monic polynomial and its unchanged squared norm by
\[
 b_n^{(\lambda)}(t)=n!P_n^{(\lambda)}(t/2;\pi/2),\qquad
 \gamma_n^{(\lambda)}=2^{1-2\lambda}n!\Gamma(n+2\lambda).
                                                               \tag{TG.11}
\]
The coefficient of \(t^n\) in (TG.9) is \(1/n!\), since
\(\arctan z=z+O(z^3)\).  Consequently \(b_n^{(\lambda)}\)
is monic.  Equation (TG.10) gives the full orthogonality identity
\[
 \int_{\mathbb R}b_n^{(\lambda)}(t)b_j^{(\lambda)}(t)
                   d\sigma_\lambda(t)
       =\gamma_n^{(\lambda)}\delta_{nj}.
                                                               \tag{TG.12}
\]
All these polynomials have real coefficients, as is also clear from
the real form of the generating function.  Differentiating (TG.9)
gives \((1+z^2)\partial_zF_\lambda=(t-2\lambda z)F_\lambda\).
Its coefficient of \(z^n\) proves
\[
 \begin{aligned}
 b_0^{(\lambda)}&=1,\qquad b_1^{(\lambda)}=t,\\
 b_{n+1}^{(\lambda)}&=t b_n^{(\lambda)}
              -n(n+2\lambda-1)b_{n-1}^{(\lambda)},\quad n\geq1,\\
 \gamma_n^{(\lambda)}/\gamma_{n-1}^{(\lambda)}
              &=n(n+2\lambda-1),\quad n\geq1.
 \end{aligned}                                                    \tag{TG.13}
\]
For the actual reference \(\lambda=1/4\), abbreviate
\(b_n=b_n^{(1/4)}\) and \(\gamma_n=\gamma_n^{(1/4)}\).  Thus
\[
 \boxed{\gamma_n=\sqrt2\,n!\Gamma(n+\tfrac12),\qquad
           \gamma_{n+1}/\gamma_n=(n+1)(n+\tfrac12).}
                                                               \tag{TG.14}
\]
These formulas agree, with the displayed Jacobian and monic factors,
with NIST DLMF \S18.19, equations (18.19.8)--(18.19.9), and \S18.22,
equation (18.22.8).

\subsection{Finite arithmetic determinants in the retained metrics}

For \(N\geq0\) let \(\mathsf\Gamma_N=\operatorname{diag}
(\gamma_0,\ldots,\gamma_N)\), and keep both matrices
\[
 \begin{aligned}
 \mathsf B_N&=\left(\int_{\mathbb R}
              b_i(t)b_j(t)|a(t)|^2d\sigma_{1/4}(t)\right)_{0\leq i,j\leq N},\\
 \mathsf J_N&=\mathsf\Gamma_N^{-1/2}\mathsf B_N
                       \mathsf\Gamma_N^{-1/2},\qquad
 D_N=\det\mathsf J_N,\quad D_{-1}=1.
 \end{aligned}                                                    \tag{TG.15}
\]
The diagonal factors are explicit invertible coordinate maps;
\(\mathsf B_N=\mathsf\Gamma_N^{1/2}\mathsf J_N
\mathsf\Gamma_N^{1/2}\) recovers every original inner product.
The integrals converge with no unproved growth estimate.  Indeed,
integration of the step function \(\lfloor x\rfloor\) gives,
initially for \(\operatorname{Re}s>1\), and then by analytic
continuation for \(\operatorname{Re}s>0\), \(s\ne1\),
\[
 \zeta(s)=\frac{s}{s-1}
         -s\int_1^\infty\{x\}x^{-s-1}\,dx.
\]
The initial equality follows by integrating each constant step on
\([n,n+1)\), or by summation by parts of the absolutely convergent
Dirichlet series.  The last integral converges absolutely and
locally uniformly in the asserted half-plane since \(0\leq\{x\}<1\).
On \(s=1/2+it\), it proves
\( |\zeta(s)|\leq1+2\sqrt{t^2+1/4}\), because
\(|s/(s-1)|=1\).  Equation (TG.2) therefore bounds \(|a(t)|\)
by a polynomial of degree three.  The exponential integrability
just proved gives all its polynomial moments.  The packet
multiplier has the same property on the tail because \(h\) is
monic, and is bounded on each compact interval by its cancelled
analytic expression in (TG.4).
Each matrix is positive definite:
the integral of the squared modulus of a nonzero polynomial is
positive, because \(a\) vanishes only at isolated real points.

\begin{theorem}[The actual theta norms as arithmetic determinants]
For all \(n\geq0\), with the initial value included,
\[
 \boxed{\mathfrak h_n=\gamma_n\frac{D_n}{D_{n-1}},\qquad
 \frac{\mathfrak h_{n+1}}{\mathfrak h_n}
   =(n+1)(n+\tfrac12)\frac{D_{n+1}D_{n-1}}{D_n^2}.}
                                                               \tag{TG.16}
\]
\end{theorem}
\begin{proof}
Let \(Q_n(t)\) be the monic polynomial orthogonal to degrees below
\(n\) in \(L^2(\mu)\).  Its existence and uniqueness follow
from the positive definite Gram matrix.  Expressing
\(Q_0,\ldots,Q_N\) in \(b_0,\ldots,b_N\) gives an upper
triangular matrix with diagonal entries one.  The determinant of
the Gram matrix is therefore unchanged, and orthogonality gives
\(\det\mathsf B_N=\prod_{j=0}^N\|Q_j\|_\mu^2\).

The algebra isomorphism
\(\Psi:\mathbb C[s]\to\mathbb C[t]\),
\((\Psi p)(t)=p(1/2+it)\), preserves the original integral
inner product.  On the affine space of degree-\(n\) monic
polynomials the map \(p\mapsto i^{-n}\Psi p\) is its monic
version and preserves the norm because \(|i^{-n}|=1\).
It takes all lower-degree polynomial subspaces onto themselves.
Thus the original monic theta residual has squared norm
\(\mathfrak h_n=\|Q_n\|_\mu^2\).  The degree-dependent phase
is used on this affine space; it is not asserted to define a linear
map on the entire polynomial algebra.

Finally \(D_N=\det\mathsf B_N/\prod_{j=0}^N\gamma_j\).
Divide this equality at \(N=n\) by the equality at \(N=n-1\),
taking an empty product to be one.  This proves the first formula.
Dividing consecutive formulas and using (TG.14) proves the second.
\end{proof}

Define \(\mathsf B_N^{(h)},\mathsf J_N^{(h)},D_N^{(h)}\)
by (TG.15) with the entire-cancelled multiplier \(a_h\) of
(TG.4).  The same proof, with no restriction on whether the packet
centres are critical, gives
\[
 \kappa_n^{(h)}=\gamma_n\frac{D_n^{(h)}}{D_{n-1}^{(h)}}.
                                                               \tag{TG.17}
\]
Combining (TG.17) and (TG.16) with (KL.24) yields the exact
packet-to-reference determinant map, for every \(m\geq0\),
\[
 \boxed{r_m^{(h)}=
  \frac{(d+m)!\,\Gamma(d+m+1/2)}{m!\,\Gamma(m+1/2)}
  \frac{D_{d+m}^{(h)}}{D_{d+m-1}^{(h)}}
  \frac{D_{m-1}}{D_m}.}
                                                               \tag{TG.18}
\]
Every determinant is strictly positive and the full degree
\(d=\sum_\rho m_\rho\) is retained.

For a packet stable under both ordinary conjugation and
\(\rho\mapsto1-\overline\rho\), and for every \(m\geq0\),
(KL.32) becomes
\[
 \boxed{\epsilon_m^2=(m+1)(m+\tfrac12)
       \frac{D_{m+1}D_{m-1}}{D_m^2}
           (r_m^{-1}-1)(1-r_{m+1}).}
                                                               \tag{TG.19}
\]
This is the exact same finite-stage allowance.  Its factor
\((m+1)(m+1/2)\) is explicit; the arithmetic determinant
quotient remains in the formula.

\subsection{Full confluent kernel and the original coordinates}

Set \(\widehat b_j(s)=i^j b_j((s-1/2)/i)\).  These are monic
polynomials in the original variable \(s\).  With
\(\mathord{\textsf{Δ}}_N=\operatorname{diag}(1,i,\ldots,i^N)\),
their actual packet Gram matrix is
\[
 \mathsf B_{N,s}^{(h)}=\mathord{\textsf{Δ}}_N^*
                  \mathsf B_N^{(h)}\mathord{\textsf{Δ}}_N.
                                                               \tag{TG.20}
\]
Indeed \(\widehat b_j(1/2+it)=i^j b_j(t)\); taking the
conjugate in the first argument gives exactly the displayed
congruence, including every phase.

Let \(\mathsf C_N\) have columns
\([\widehat b_j]_h\in E_h=\mathbb C[s]/(h)\), in the original
monomial remainder coordinates.  For \(N\geq d-1\),
\[
 \boxed{K_N=\mathsf C_N(\mathsf B_{N,s}^{(h)})^{-1}
                           \mathsf C_N^*,\qquad
  \mathbf G_{N-d}=U_{\varepsilon_h}^*K_N^{-1}
                           U_{\varepsilon_h}.}
                                                               \tag{TG.21}
\]
Here \(\mathbf G_{-1}\) is the fixed initial-section metric,
and \(\varepsilon_h=(j_h(g/h))^{-1}\) is the same full local
unit as before.  To prove the first identity, choose the actual
monic orthogonal polynomial basis through degree \(N\).  If
\(T\) is its invertible coefficient matrix in the
\(\widehat b_j\) basis, then
\[
 T^*\mathsf B_{N,s}^{(h)}T
     =\operatorname{diag}(\kappa_0,\ldots,\kappa_N).
\]
Inverting gives
\((\mathsf B_{N,s}^{(h)})^{-1}
 =T\operatorname{diag}(\kappa_j^{-1})T^*\).
The columns of \(\mathsf C_NT\) are the full remainders
\([p_j]_h\).  Substitution proves exactly
\(K_N=\sum_{j=0}^N[p_j]_h[p_j]_h^*/\kappa_j\), the original
interpolation kernel.  The first \(d\) monic polynomials span all
remainders, so \(\mathsf C_N\) is onto and \(K_N\) is positive
definite.  The second identity is the fixed-unit minimum-lift
formula already proved in (KL.27), with the exact index
\(N=d+m\).  No single remainder column is required to be nonzero
or invertible.

For a centre \(\rho\), put \(t_\rho=(\rho-1/2)/i\) and
write \(f^{[\ell]}=f^{(\ell)}/\ell!\).  The full Taylor
coordinates of these columns are explicitly
\[
 \widehat b_j^{[\ell]}(\rho)
     =i^{j-\ell}b_j^{[\ell]}(t_\rho),
                 \qquad 0\leq\ell<m_\rho.
                                                               \tag{TG.22}
\]
This follows by differentiating the affine argument \(\ell\)
times.  CRT converts these full Taylor coordinates to the monomial
remainder coordinates, and multiplication by
\(U_{\upsilon_h}\) retains all derivatives of \(g/h\).  Thus
(TG.21) computes the entire confluent kernel, including nilpotent
orders and critical centres.

\subsection{The finite map to Romik's gamma parameter}

Romik's expansion uses \(P_n^{(3/4)}(t/2;\pi/2)\), with the
unchanged variable \(t/2\); see his \S3 and Appendix A.2.
The exact finite change to the parameter in (TG.3) is
\[
 \begin{aligned}
 b_n^{(1/4)}(t)
   &=\sum_{k=0}^{\lfloor n/2\rfloor}
       \frac{n!}{(n-2k)!}\binom{1/2}{k}
                       b_{n-2k}^{(3/4)}(t),\\
 b_n^{(3/4)}(t)
   &=\sum_{k=0}^{\lfloor n/2\rfloor}
       \frac{n!}{(n-2k)!}\binom{-1/2}{k}
                       b_{n-2k}^{(1/4)}(t).
 \end{aligned}                                                    \tag{TG.23}
\]
To prove these formulas, divide the two generating functions:
\(F_{1/4}(z,t)=(1+z^2)^{1/2}F_{3/4}(z,t)\).
The binomial series at zero and coefficient comparison prove the
first formula.  Multiplication by \((1+z^2)^{-1/2}\) proves the
inverse.  Each finite coefficient matrix is triangular with diagonal
one, preserves parity, and is invertible.  Applied to the same
measure, it transports Gram matrices by its full congruence and
transports every Taylor column by the same coefficient matrix.
Consequently the kernel expression (TG.21) is invariant under this
change of polynomial basis.

The two reference measures themselves have the exact Hilbert-space
map
\[
 L^2(\sigma_{1/4})\longrightarrow L^2(\sigma_{3/4}),
 \qquad F(t)\longmapsto
        \frac{\Gamma(1/4+it/2)}{\Gamma(3/4+it/2)}F(t).
                                                               \tag{TG.24}
\]
The gamma quotient has neither zeros nor poles for real \(t\).
Cancellation of its squared modulus proves the isometry and its
inverse; this isometry has a multiplier, whereas (TG.23) is the
finite polynomial coordinate map.  Their respective domains and
metrics have both been given explicitly.

Romik's coefficient asymptotics concern the expansion of \(\Xi\)
in a fixed gamma-orthogonal basis.  Inoue's \S2 similarly treats
expansion coefficients and convergence.  In the present construction
the arithmetic norms are the determinants (TG.16).  Equations
(TG.20)--(TG.24) are the exact maps connecting those polynomial
presentations to these determinants; no assertion that their
coefficient asymptotics already estimate \(D_n\) is used.

\subsection{Exact scope of the exponential-weight theorem}

The original theorem of Lubinsky--Mhaskar--Saff does include tail
order one.  In their 1988 paper, Definition 2.1 requires
\(W=e^{-Q}\), with \(Q\) even and continuous, \(Q'\) existing
on \((0,\infty)\), \(xQ'(x)\) bounded at zero, and, eventually,
\[
 Q'>0,\qquad x^2|Q'''|/Q'\leq C,\qquad
            1+xQ''/Q'\longrightarrow\alpha>0.
                                                               \tag{TG.25}
\]
Theorem 2.3 admits a finite generalized Jacobi factor, a lower-degree
polynomial exponential, and a bounded nonnegative factor
\(\Psi(t)\to1\).  Its Theorem 2.4 replaces the polynomial
exponential by \(V\), but requires
\(\log V(t)/Q(t)\to0\), as well as bounded reciprocal
approximation on the expanding interval by polynomials of degree
\(o(n)\).  The precise statements are on printed pages 67--70;
their proof is on pages 80--81.  These hypotheses concern the
full weight, not its gamma factor alone.

We prove a concrete obstruction to inserting the actual arithmetic
multiplier into those displayed hypotheses.  Hardy's theorem
supplies real numbers \(T_j\to\infty\) with
\(g(1/2+iT_j)=0\).  Removing a fixed full polynomial packet
leaves an infinite subsequence of these zeros.  At each remaining
\(T_j\), the multiplier \(a_h\) of (TG.4) has the local form
\[
 a_h(t)=(t-T_j)^{\ell_j}v_j(t),\qquad
 \ell_j\geq1,\quad v_j(T_j)\ne0,
                                                               \tag{TG.26}
\]
where \(v_j\) is analytic in a neighbourhood of \(T_j\).
This follows from the Taylor series at an isolated zero, with its
entire order \(\ell_j\) retained.

Let \(B(t)>0\) be any continuous comparison amplitude near all
sufficiently large \(T_j\), and suppose it has no zeros there.
For every \(j\) and every \(\eta>0\), continuity and (TG.26)
give a nonempty interval \(I_{j,\eta}\) about \(T_j\) on which
\[
                 |a_h(t)|/B(t)<\eta.
                                                               \tag{TG.27}
\]
Thus a decomposition with this ratio as the bounded multiplier
\(\Psi\) cannot satisfy \(\Psi\to1\).  The obstruction
holds on intervals of positive measure, so changing the value only
at each zero does not repair it.  A finite generalized Jacobi
factor has only finitely many real zeros or singularities; choosing
the subsequence beyond them leaves (TG.27) unchanged.

For the logarithmic hypothesis choose any eventually positive
continuous \(Q\).  At a remaining \(T_j\), shrink the
neighbourhood so that \(Q(t)\leq Q(T_j)+1\), and choose
\(\eta=\exp(-j(Q(T_j)+1))\) in (TG.27).  On the resulting
positive-measure interval,
\[
 \frac{\log(|a_h(t)|/B(t))}{Q(t)}<-j
                  \quad(t\ne T_j).
                                                               \tag{TG.28}
\]
Hence this ratio cannot be \(V\) with \(\log V/Q\to0\).
An additional factor \(\Psi\to1\) is bounded away from zero
and infinity on the tail and does not remove the contradiction.
This disproves the indicated application of Theorems 2.3--2.4,
including almost-everywhere modifications of the weight.
It makes no claim that other asymptotic methods are unavailable.
The exact positive Gram matrices (TG.15)--(TG.21) retain the
arithmetic information at all those zeros and are available for
such further analysis.

\subsection{Primary source locators}

The content-reading route and source hashes are recorded in
\path{work/theta_norm_literature_route_20260912.json} and
\path{work/theta_norm_literature_sources/download_receipt.json}.
The principal public locators are:
\begin{itemize}
 \item D.~S.~Lubinsky, H.~N.~Mhaskar and E.~B.~Saff,
 \emph{A proof of Freud's conjecture for exponential weights},
 Constructive Approximation 4 (1988), 65--83, Definition 2.1,
 Theorems 2.3--2.5 and \S4:
 \url{https://math.vanderbilt.edu/saffeb/texts/81.pdf}.
 \item Dan Romik, \emph{Orthogonal polynomial expansions for the
 Riemann xi function in the Hermite, Meixner--Pollaczek, and
 continuous Hahn bases}, \S3.1--3.3 and Appendix A.2:
 \url{https://www.math.ucdavis.edu/~romik/data/uploads/papers/riemannxi-final.pdf}.
 \item Hiroto Inoue, \emph{Expansion of Dirichlet L-function on the
 critical line in Meixner--Pollaczek polynomials}, arXiv:1412.1220,
 \S2.1--2.2:
 \url{https://arxiv.org/abs/1412.1220}.
 \item NIST DLMF, equations (18.19.8)--(18.19.9), (18.22.8):
 \url{https://dlmf.nist.gov/18.19},
 \url{https://dlmf.nist.gov/18.22}.
 \item G.~H.~Hardy, \emph{Sur les z\'eros de la fonction
 \(\zeta(s)\) de Riemann}, Comptes rendus 158 (1914),
 1012--1014, \S1--3.  The exact page images and transcriptions
 are linked in the machine-readable source route.
\end{itemize}
